%% notes-tex/025-additive-baselines/sections/2-data.tex

\section{Data and the two splits\secstatus{todo}}
\label{sec:data}

\notesec{The label}
Write $f_i$, $f_{ij}$ and $f_{ijk}$ for the fitness of the single, double and triple
mutant relative to wild type, and $\varepsilon_{ij} = f_{ij} - f_i f_j$ for the
digenic interaction. The trigenic interaction score is
%
\begin{equation}
  \tau_{ijk} \;=\; f_{ijk} - f_i f_j f_k
  - f_i\,\varepsilon_{jk} - f_j\,\varepsilon_{ik} - f_k\,\varepsilon_{ij} ,
  \label{eq:tau}
\end{equation}
%
\external{Kuzmin et al.\ 2018, \texttt{10.1126/science.aao1729}; the definition of
the released adjusted score, not re-derived here}
a third-order residual: the multiplicative part of fitness and every pairwise term
are subtracted before the number is reported, so $\tau$ is by construction the part
of the triple phenotype that no single or double mutant predicts. The loaders store
it as \texttt{gene\_interaction}.

\notesec{The records}
Subset S0 of the 025 build is its 376{,}732 trigenic records over 4{,}352 genes.
Matching them to the 010 build by genotype recovers every one; 376{,}526 labels are
bit-identical and the remaining 206 differ by at most $5.6 \times 10^{-17}$
(\file{experiments/025-solid-growth/scripts/label_parity_010_vs_025.py}). Every
model below is fit on these records and nothing else.

\notesec{The screen structure}
The records contain 739{,}315 distinct unordered gene pairs, of which 420 occur five
or more times. Those 420 account for 376{,}733 pair instances, one more than the
record count: every record contains exactly one recurring pair, and a single record
contains two. The 420 pairs are the Kuzmin query doubles. Grouping records by their
recurring pair recovers the screen exactly, and any model with a per-pair offset can
absorb whatever is common to a screen.

\notesec{The two arms}
Arm R is 010's own split, random over records, carried into the 025 index space by
genotype identity. Arm Q assigns whole query pairs to parts, so no screen appears in
more than one part. The two arms hold out nearly the same number of records and
differ only in what is held out.

%% SOURCE: experiments/025-solid-growth/scripts/additive_baselines_025_panels.py ->
%% tables/t1-arms.tex, from the two split artifacts named in the table.
\begin{table}[htbp]
\centering
\footnotesize
\caption{The two arms. Both partition the same 376{,}732 records. Each part of arm
R carries every one of the 420 query pairs; the parts of arm Q carry disjoint sets
of them. The split artifacts are the files the trainer loads, under
\texttt{experiments/025-solid-growth/results/}.}
\label{tab:arms}
%% GENERATED by experiments/025-solid-growth/scripts/additive_baselines_025_panels.py
%% SOURCE: the result files named in that script's docstring. Do not edit by hand.
\begin{tabular}{lll}
\toprule
 & Arm R & Arm Q \\
\midrule
Held-out unit & record & query pair \\
Recurring query pairs & 420, shared by every part & 420, no pair in two parts \\
Query pairs, train / val / test &  & 331 / 43 / 46 \\
Records, train / val / test & 301{,}386 / 37{,}673 / 37{,}673 & 301{,}236 / 37{,}705 / 37{,}791 \\
Split artifact & \file{pinned_splits_from_010_seed_42.json.gz} & \file{query_pair_disjoint_splits_025.json.gz} \\
Transformer run & GH 1598 & GH 1640, and Table~\ref{tab:disjointruns} \\
\bottomrule
\end{tabular}

\end{table}

\notesec{What arm Q holds out}
The query pair, not the genes. Every array gene of the validation and test parts
occurs in training, as do 78 of the 86 validation and 88 of the 92 test query-pair
genes, so 81 percent of validation records and 91 percent of test records have all
three genes in the training set (Table~\ref{tab:armqgenes}). A model on arm Q is
therefore asked to predict an unseen pair of mostly seen genes with a seen third
gene, which is the generalization that scoring an unmeasured triple requires. The
test part spreads its records over 1{,}182 distinct array genes against the
validation part's 4{,}003; the two parts are different kinds of screen sample, and
no analysis here uses that difference.

%% SOURCE: experiments/025-solid-growth/scripts/additive_baselines_025_panels.py ->
%% tables/t4-armq-genes.tex, from the arm Q split artifact and the per-triple gene
%% table of the 025 build.
\begin{table}[htbp]
\centering
\footnotesize
\caption{What the held-out parts of arm Q share with its training part. A query pair
is the most frequent recurring pair a record carries, the rule that built the split;
the array gene is the record's third gene.}
\label{tab:armqgenes}
%% GENERATED by experiments/025-solid-growth/scripts/additive_baselines_025_panels.py
%% SOURCE: the result files named in that script's docstring. Do not edit by hand.
\begin{tabular}{lrr}
\toprule
 & Validation & Test \\
\midrule
Records & 37{,}705 & 37{,}791 \\
Query pairs held out & 43 & 46 \\
Query pairs also in training & 0 & 0 \\
Distinct array genes & 4{,}003 & 1{,}182 \\
Array genes present in training & 4003 of 4003 & 1182 of 1182 \\
Query-pair genes present in training & 78 of 86 & 88 of 92 \\
Records with all three genes in training & 81.3\% & 91.4\% \\
\bottomrule
\end{tabular}

\end{table}
